\section{Multi-Agent: ganancia de colaboración y costo de comunicación}

\subsection{Cuándo usar múltiples agentes}
El curso no fomenta ``paralelizar por paralelizar''. Multi-agent es adecuado para escenarios donde la tarea es descomponible, los límites de los roles son claros y el protocolo de comunicación es estable. Si la tarea tiene un acoplamiento fuerte o el estado intermedio se comparte en gran medida, un Agent monolítico con memory fuerte suele ser más estable.

\begin{knowledgebox}{Origen de la ganancia en sistemas multiagente}
\begin{itemize}
    \item la especialización de roles reduce la complejidad de la estrategia del agente monolítico;
    \item la exploración en paralelo aumenta la cobertura de soluciones candidatas;
    \item la revisión cruzada reduce el riesgo de alucinación de una sola ruta.
\end{itemize}
\end{knowledgebox}

\begin{table}[H]
\centering
\begin{tabular}{p{3.3cm}p{3.8cm}p{3.9cm}}
\toprule
\textbf{Dimensión} & \textbf{Agent monolítico} & \textbf{Multi-Agent} \\
\midrule
Gestión del contexto & Simple, pero puede congestionarse & Distribuida, pero requiere protocolo de sincronización \\
Eficiencia de ejecución & Secuencial, fácil de depurar & Alto potencial de paralelismo, coordinación compleja \\
Localización de errores & La localización centralizada es más intuitiva & Es necesario distinguir errores de rol y errores de comunicación \\
Estructura de costos & Relativamente estable y predecible & Alto costo de comunicación e inferencia repetida \\
\bottomrule
\end{tabular}
\caption{Disyuntivas entre arquitectura monolítica y multiagente}
\end{table}

\begin{importantbox}{Principios de implementación de sistemas multiagente}
Primero se diseña el protocolo y después se agregan roles. El protocolo debe incluir al menos: formato de los mensajes, alcance del estado compartido, reglas de resolución de conflictos, y una estrategia de tiempo de espera y reversión.
\end{importantbox}

\begin{warningbox}{Modos de falla más comunes en sistemas multiagente}
\begin{itemize}
    \item la superposición de responsabilidades entre roles produce trabajo duplicado;
    \item la información de comunicación se infla y consume el presupuesto de contexto;
    \item la falta de un validador unificado da lugar a que «se confirmen entre sí pero el resultado global sea incorrecto».
\end{itemize}
\end{warningbox}

\subsection{Resumen de la sección}
Multi-agent no es más fuerte por defecto, sino condicionalmente más fuerte. Lo que importa es la calidad de la descomposición de la tarea y el diseño del protocolo de comunicación, no la cantidad de agents en sí.

